\begin{frame}{Konstrukcija pravokotnice}
\begin{columns}
	\begin{column}{0.5\textwidth}
		\begin{itemize}[<+->]
			\item Dana sta premica \(p\) in točka \(T\).
			\item Nariši lok \(k\) s središčem v \(T\).
			\item Lok \(k\) seče premico \(p\) v točkah \(A\) in \(B\).
			\item Nariši lok \(m\) s središčem v \(A\).
			\item Nariši lok \(n\) s središčem v \(B\) in z enakim polmerom.
			\item Loka se sečeta v točki \(C\).
			\item Premica \(TC\) je pravokotna na \(p\).
		\end{itemize}
	\end{column}
	\begin{column}{0.45\textwidth}
		\begin{center}
			\includegraphics<1>[width=50mm]{./slike/fig-1.png}%
			\includegraphics<2>[width=50mm]{./slike/fig-2.png}%
			\includegraphics<3>[width=50mm]{./slike/fig-3.png}%
			\includegraphics<4>[width=50mm]{./slike/fig-4.png}%
			\includegraphics<5>[width=50mm]{./slike/fig-5.png}%
			\includegraphics<6>[width=50mm]{./slike/fig-6.png}%
			\includegraphics<7>[width=50mm]{./slike/fig-7.png}%
		\end{center}
  	\end{column}
\end{columns}
\end{frame}
